\documentclass[11pt]{article}
\usepackage{mathblatt}
\makeatletter
\setlength{\mbpfbe}{0mm}% keine Punkte (Bauregel 6.4)
\makeatother
% Rechenraum als Karo (Bauregel 6.9): n Rechenschritte = n mal 10 mm
\newcommand{\karo}[1]{\par\smallskip\noindent\begin{tikzpicture}
  \clip (0,0) rectangle ({\linewidth-1mm},{#1*10mm});
  \draw[black!16,line width=0.3pt,step=5mm] (0,0) grid ({\linewidth},{#1*10mm});
  \end{tikzpicture}\par}
\newcommand{\linie}[1]{\,{\color{black!45}\rule[-1pt]{#1}{0.4pt}}\,}
\newcommand{\pkt}{\textperiodcentered{}}
\begin{document}

\pfheftstil
\pflaufkopf[P10 \pkt{} mehr]{Prozent \pkt{} Lösungen}{PS10-2}\pfseitevon
\pfheftkopf{Prozent}{Lösungen \pkt{} mehr}
\begin{pfloesungb}{}
\lzb{1.}{}{Spannweite 2,8 Mio.; Mittel 8,4 Mio.; Erhöhung $\approx$ 14,7 \%}{Spannweite 2,8 Mio.; Mittel 8,4 Mio.; Erhöhung $\approx$ 14,7 \%}
\end{pfloesungb}
\begin{pfloesungb}{}
\lzb{2.}{}{C: 7 von 100 Bäumen wachsen nicht an (dritte Option)}{7 von 100 Bäumen wachsen nicht an (dritte Option)}
\end{pfloesungb}
\begin{pfloesungb}{}
\lzb{3.}{}{7,5 \%}{7,5 \%}
\end{pfloesungb}
\begin{pfloesungb}{}
\lzb{4.}{}{18 €}{18 €}
\end{pfloesungb}
\begin{pfloesungb}{}
\lzb{5.}{}{2,7 \%}{2,7 \%}
\end{pfloesungb}
\begin{pfloesungb}{}
\lzb{6.}{}{300 € · 0,03 = 9,00 €}{300 € · 0,03 = 9,00 €}
\end{pfloesungb}
\begin{pfloesungb}{}
\lzb{7.}{}{Spannweite 240 $-$ 120 = 120; Mittwoch 180 + 60 = Sonntag}{Spannweite 240 $-$ 120 = 120; Mittwoch 180 + 60 = 240 = Sonntag}
\end{pfloesungb}
\begin{pfloesungb}{}
\lzb{8.}{}{nein: Studie 1 $\approx$ 76,8; Studie 2 = 72,0 Tausend Liter}{nein: Studie 1 $\approx$ 76,8; Studie 2 = 72,0 Tausend Liter}
\end{pfloesungb}
\begin{pfloesungb}{}
\lzb{9.}{}{Aussage 1 falsch; Etwa zwei Drittel (69 \%) der Mädchen nutzen einen Fernseher (vier Fünftel gilt für die Jungen)}{Aussage 1 falsch; Etwa zwei Drittel (69 \%) der Mädchen nutzen einen Fernseher (vier Fünftel gilt für die Jungen)}
\end{pfloesungb}
\begin{pfloesungb}{}
\lzb{10.}{}{35,1 \% + 11,8 \% + 4,4 \% = 51,3 \%}{35,1 \% + 11,8 \% + 4,4 \% = 51,3 \%}
\end{pfloesungb}
\begin{pfloesungb}{}
\lzb{11.}{}{90 : 450 = 20 \%}{90 : 450 = 20 \%}
\end{pfloesungb}
\begin{pfloesungb}{}
\lzb{12.}{}{11 : 87 $\approx$ 12,6 \%; Sektor mit $\approx$ 46° ab dem vorgegebenen Radius;  $\approx$ 46°}{11 : 87 $\approx$ 12,6 \%; Sektor mit $\approx$ 46° ab dem vorgegebenen Radius; $\approx$ 46°}
\end{pfloesungb}
\begin{pfloesungb}{}
\lzb{13.}{}{Guthaben 14. Geburtstag: 927,27 € + 27,82 € + 300 € = 1255,09 €; Zinsen 15. Geburtstag: 1255,09 € · 0,03 = 37,65 €}{Guthaben 14. Geburtstag: 927,27 € + 27,82 € + 300 € = 1255,09 €; Zinsen 15. Geburtstag: 1255,09 € · 0,03 = 37,65 €}
\end{pfloesungb}
\begin{pfloesungb}{}
\lzb{14.}{}{1500 € · 1,02$^{4}$ $\approx$ 1623,65 €}{1500 € · 1,02$^{4}$ $\approx$ 1623,65 €}
\end{pfloesungb}
\end{document}